%% Chapitre 13 — Intégrer Claude Code dans ton workflow quotidien.
%% RÈGLE : uniquement des commandes sémantiques bookforge. Aucun terme
%% du glossaire n'est réintroduit ici (tous déjà vus dans les chapitres
%% précédents) : aucun \terme{} n'est marqué.

\chapitre{Intégrer Claude Code dans ton workflow quotidien}

\introChapitre{
  Dernier chapitre : on assemble tout en habitudes durables. Comment l'agent
  prend sa place dans ta journée, ton dépôt et ton équipe, sans devenir une
  béquille ni un risque. On clôt sur une routine réaliste et sur la manière de
  continuer à progresser quand l'outil évoluera.
}

\paragraphe{Tu sais maintenant l'essentiel : reformuler une tâche pour la rendre
délégable, soigner le contexte que tu donnes, relire le diff au lieu de faire
confiance aveuglément, poser des garde-fous via les permissions, et orchestrer
plusieurs subagents quand le chantier le mérite. Ces cinq réflexes ne valent que
s'ils deviennent des \emphase{habitudes}. Ce chapitre les transforme en routine.}

\section{Une journée type : où l'agent aide, où tu gardes la main}

\paragraphe{Le matin, tu ouvres une session sur une intention floue : « comprends
pourquoi le paiement échoue en recette ». L'agent excelle là. Il lit le code,
suit les appels, te résume le chemin du bug. Tu n'aurais pas tapé ces dix
\code{grep} toi-même. Tu \emphase{gardes la main} sur le diagnostic final : c'est
toi qui décides si l'hypothèse tient.}

\paragraphe{En milieu de matinée, le travail répétitif. Migrer trente fichiers
vers une nouvelle signature, écrire les tests d'un module déjà spécifié, mettre
à jour la doc après un renommage. Tâches bornées, critère de succès clair :
délègue sans hésiter, relis le diff par paquets.}

\paragraphe{L'après-midi, une décision d'architecture. Là, l'agent propose mais
ne tranche pas. Demande-lui deux options et leurs compromis, puis \emphase{décide
toi-même}. Pareil pour tout ce qui touche au sensible : secrets, migration de
base en production, suppression massive. L'agent reste un exécutant brillant et
faillible ; l'arbitrage et la responsabilité restent à toi.}

\begin{note}
Une bonne répartition tient en une phrase : délègue ce que tu saurais vérifier,
garde ce que tu serais incapable de relire. Si tu ne peux pas juger le résultat,
tu ne délègues pas, tu paries.
\end{note}

\section{Conventions de dépôt et d'équipe}

\paragraphe{Dès que vous êtes plusieurs sur un dépôt, l'usage de l'agent doit
être \emphase{partagé}, pas réinventé par chacun. Le pivot, c'est la mémoire
projet : un fichier versionné qui dit à l'agent comment se comporter ici. Tout le
monde en hérite, tout le monde l'améliore.}

\paragraphe{Un squelette minimal suffit pour démarrer une équipe :}

\begin{liste}
  \element{Conventions de code et de tests : commandes de lint et de build,
  style attendu, ce qui est interdit (toucher aux migrations, par exemple).}
  \element{Permissions par défaut versionnées : ce qui s'exécute seul, ce qui
  demande approbation. Les garde-fous deviennent un choix d'équipe, pas une
  préférence individuelle.}
  \element{Slash commands partagées pour les tâches récurrentes du projet :
  « prépare la revue », « écris les tests du module X » — une fois écrites, tout
  le monde les rejoue.}
\end{liste}

\paragraphe{Côté collaboration humaine, garde une règle simple : un changement
produit par l'agent passe la \emphase{même revue} qu'un changement humain. Ni
plus laxiste parce que « c'est l'agent », ni plus méfiant. C'est du code à
relire, point. Indique dans le message de commit qu'il a été assisté : ça aide
l'équipe à calibrer sa vigilance.}

\begin{avertissement}
Ne committe jamais la sortie d'une session que personne n'a relue, surtout sur
une branche partagée. « Ça compilait » n'est pas une revue. Un diff vert peut
cacher une hallucination parfaitement plausible.
\end{avertissement}

\section{Construire tes propres routines}

\paragraphe{Les bonnes habitudes ne s'improvisent pas à chaud, elles
s'\emphase{outillent}. À chaque fois que tu te surprends à réécrire la même
consigne — « relis d'abord, ne touche pas aux tests, montre-moi un plan avant
d'agir » — tu tiens une routine à figer.}

\paragraphe{Deux familles de routines reviennent tout le temps. Les routines de
délégation : un prompt-type qui pose le but, le périmètre et le critère de
succès, prêt à remplir. Les routines de vérification : ta liste de contrôle
avant d'accepter un diff.}

\begin{extraitcode}{text}
Avant d'accepter un diff de l'agent :
- [ ] J'ai lu chaque ligne, pas seulement le résumé.
- [ ] Les tests passent ET un test échoue si je casse le code volontairement.
- [ ] Aucune fonction/API inventée : j'ai vérifié les noms douteux.
- [ ] Le périmètre est respecté : rien n'a changé hors de la demande.
- [ ] Secrets, migrations, prod : intacts ou validés explicitement.
\end{extraitcode}

\paragraphe{Encapsule ces routines là où elles vivront : une slash command pour
la délégation, un hook pour automatiser le formatage et les tests après chaque
modification. Ta liste de vérification, elle, reste sous \emphase{tes} yeux —
c'est ton jugement, pas celui de l'agent, qui valide.}

\section{Mesurer si l'agent te fait vraiment gagner}

\paragraphe{Le danger silencieux, c'est l'illusion de productivité : tu passes
vingt minutes à corriger un diff que tu aurais écrit en quinze. Pour le voir, il
faut \emphase{mesurer}, même grossièrement.}

\paragraphe{Garde en tête trois signaux, sur quelques semaines :}

\begin{liste}
  \element{Temps réel de bout en bout : briefer plus relire plus corriger,
  comparé à « je l'aurais fait moi-même ». Sur les tâches répétitives bornées,
  l'agent gagne presque toujours ; sur le pointu très contextuel, souvent non.}
  \element{Taux de reprise : combien de fois tu relances ou réécris derrière lui.
  S'il est élevé sur un type de tâche, le problème est ta consigne ou le
  périmètre, pas l'agent.}
  \element{Coût en jetons rapporté à la valeur : une session qui brûle beaucoup
  pour un résultat que tu jettes n'est pas un gain, même si elle « marche ».}
\end{liste}

\paragraphe{La conclusion utile n'est pas « l'agent est bon ou mauvais », mais
« sur \emphase{quelles} tâches il me fait gagner ». Délègue franchement celles-là,
reprends la main sur les autres. Cette carte personnelle vaut plus que n'importe
quel conseil général.}

\section{Rester à jour quand l'outil change}

\paragraphe{Claude Code évolue vite : commandes, options, tarifs, écosystème
d'outils branchés en MCP. Les détails de ce livre dateront ; son
\emphase{esprit}, non. C'est voulu.}

\paragraphe{Sépare donc deux niveaux. Les \emphase{détails} — nom exact d'un
drapeau, prix au jeton, dernière fonctionnalité — se vérifient à la source, dans
la documentation officielle, au moment où tu en as besoin. Ne les mémorise pas :
ils périment. Les \emphase{principes} — penser la délégation, soigner le contexte,
vérifier le diff, encadrer par les permissions, orchestrer — ne bougent pas. Ils
survivront à trois changements d'interface.}

\begin{astuce}
Quand une nouveauté apparaît, ne demande pas « comment ça marche » mais « quel
principe connu ça sert ». Un nouvel outil branché en MCP, c'est juste plus de
contexte ou plus d'actions : tu sais déjà l'évaluer, le borner, le vérifier.
\end{astuce}

\section{Tes prochaines semaines avec Claude Code}

\paragraphe{Termine sur du concret. Voici un plan de montée en compétence sur un
mois, une étape par semaine, à adapter à ton rythme.}

\begin{listeNumerotee}
  \element{\emphase{Semaine 1 — déléguer pour de vrai.} Confie à l'agent trois
  tâches bornées par jour : un test, un petit refactor, une correction. Relis
  chaque diff intégralement. Objectif : ressentir où il aide.}
  \element{\emphase{Semaine 2 — encadrer.} Écris la mémoire projet de ton dépôt
  principal et fixe tes permissions par défaut. Pose un hook qui lance tes tests
  après chaque modification. Tu construis tes garde-fous.}
  \element{\emphase{Semaine 3 — industrialiser.} Transforme tes deux consignes
  les plus répétées en slash commands. Tiens ta liste de vérification à chaque
  diff jusqu'à ce qu'elle devienne un réflexe.}
  \element{\emphase{Semaine 4 — orchestrer et mesurer.} Découpe un vrai chantier
  en sous-tâches confiées à des subagents. En parallèle, note tes trois signaux
  pour dresser ta carte « où l'agent me fait gagner ».}
\end{listeNumerotee}

\paragraphe{Après ce mois, tu n'auras plus besoin de plan : tu auras des
habitudes. C'est tout l'enjeu de ce livre — non pas savoir ce que fait l'outil,
mais savoir \emphase{travailler avec}, lucidement, sans naïveté ni méfiance
paralysante.}

\begin{astuce}
Tu sais coder, et tu sais maintenant déléguer, vérifier, encadrer, orchestrer.
Ce ne sont pas des recettes pour un outil donné : c'est une façon de penser ton
métier qui restera vraie quand l'agent d'aujourd'hui aura été remplacé par celui
de demain. Ouvre une session, confie-lui ta prochaine tâche, et relis son diff
avec ces yeux-là. C'est là que ton vrai apprentissage commence.
\end{astuce}

\resumeChapitre{
  \element{\important{Répartition journalière} : délègue l'exploration et le mécanique répétitif, garde le diagnostic final, les décisions d'architecture et tout ce qui touche au sensible.}
  \element{\important{Mémoire projet d'équipe} : versionne conventions, permissions et slash commands dans le dépôt — un bien commun que tout le monde améliore et dont tout le monde hérite.}
  \element{\important{Routines outillées} : encapsule tes prompts-types en slash commands, automatise le formatage et les tests en hooks, tiens ta liste de vérification sous les yeux.}
  \element{\important{Mesurer le gain réel} : note temps total, taux de reprise et coût en jetons sur quelques semaines pour dresser ta carte personnelle « où l'agent me fait gagner ».}
  \element{\important{Principes durables, détails périssables} : les fonctionnalités et les prix changent ; déléguer, vérifier, encadrer, orchestrer restent vrais quel que soit l'outil de demain.}
}
